

%قبل از شروع مسئله ترجیحاً چیزی ننویسید. اما اگر خواستید بنویسید، در قالب لیستی (فقرات و شمارش) بنویسید.
%برای اضافه کردن نمره، بعدی از شروع‌مسئله، [x] را اضافه کنید و به جای x  نمره مورد نظر را قرار دهید. برای زیرمسئله‌ها هم به همین صورت است.
%\شروع‌مسئله[10]
\شروع‌مسئله []
%\متن‌سیاه{سوال عملی اول:}\\
%\متن‌سیاه{دراپ‌تست}\\
اگر از بالای ساختمانی با ارتفاع (\متن‌لاتین{H}) به سمت راست لیزر زده شود، برای هر کدام از ساختمان‌های با ارتفاع بیشتر از (\متن‌لاتین{H}) مانند (\متن‌لاتین{h})، به اندازه‌ی (\متن‌لاتین{h - H}) خرابی به بار می ‌آید. اگر جمع این مقدارها را برای هر ساختمان بدست بیاوریم حاصل برابر خواهد شد با : (\متن‌لاتین{S - c * H}) که (\متن‌لاتین{S}) جمع ارتفاع ساختمان‌های بلندتر از (\متن‌لاتین{H}) است که در سمت راست ساختمان با ارتفاع (\متن‌لاتین{H}) قرار دارند. 
\\بنابر این در ابتدا با درخت \متن‌لاتین{AVL} دو  مقدار را برای هر کدام از ساختمان‌ها بدست می ‌آوریم :‌ جمع ارتفاع ساختمان‌های بلندتر که در سمت راست این ساختمان قرار دارند(\متن‌لاتین{sum[i]}) و تعداد ساختمان‌های بلند‌تر با این ویژگی (\متن‌لاتین{count[i]}).
\\برای این کار یک درخت \متن‌لاتین{AVL} داریم که در ابتدا خالی است و به ترتیب از راست به چپ اعداد را به این درخت اضافه می‌کنیم و در طول عملیات \متن‌لاتین{insert} دو مورد گفته شده را هم بدست می‌آوریم. به این شکل که برای هر راس در این درخت دو مقدار را نگه می‌داریم: اول جمع اعداد زیر درخت آن (\متن‌لاتین{S[i]}) و دوم تعداد رئوس زیر درخت آن راس (\متن‌لاتین{C[i]}). حال در عملیات \متن‌لاتین{insert} یک راس جدید از ریشه شروع می‌کنیم و اگر از راسی که روی آن هستیم بزرگتر بود به فرزند راست می‌رویم و در غیر این صورت به فرزند چپ (همانطور که در \متن‌لاتین{insert} عادی عمل می‌کنیم). در هر مرحله اگر روی راس \متن‌لاتین{i} باشیم و بخواهیم راس \متن‌لاتین{j} را به درخت اضافه کنیم، اگر به فرزند چپ \متن‌لاتین{i} رفتیم یعنی \متن‌لاتین{h[i] > h[j]}، در این صورت تمام زیردرخت فرزند راست راس \متن‌لاتین{i}  از راس \متن‌لاتین{j} مقدار بیشتری دارند و چون از راست به چپ مقادیر ارتفاع ساختمان‌ها را وارد کردیم همه‌‌ی آن‌ها سمت راست ساختمان \متن‌لاتین{j} قرار دارند. پس \متن‌لاتین{sum[j]} را با \متن‌لاتین{s[i.right]} و \متن‌لاتین{count[j]} را با \متن‌لاتین{c[i.right]} جمع می‌زنیم (به علاوه‌ی مقدار خود راس \متن‌لاتین{i}). اگر هم به فرزند راست \متن‌لاتین{i} رفتیم کاری نمیکنیم. در پایان عملیات \متن‌لاتین{insert} هم درخت را متوازن می‌کنیم و در طی آن مقادیر \متن‌لاتین{c[i]} و \متن‌لاتین{s[i]} را برای رئوسی که زیردرختشان تغییر کرده \متن‌لاتین{update} می‌کنیم (از برگ ها شروع کرده و هرکس جمع دو فرزندش + مقدار خودش می‌شود).
\\حال باید برای هر ساختمان اولین ساختمان بلندتر یا هم ارتفاع سمت راست آن را بدست بیاوریم (\متن‌لاتین{r[i]}). برای این منظور یک پشته‌ی صعودی در نظر می‌گیریم که در هر لحظه ساختمان‌هایی که می‌توانند کاندید اولین ساختمان بلندتر سمت راست ساختمانی باشند در آن قرار دارند. در ابتدا پشته خالی است. از سمت راست شروع می‌کنیم و به هر ساختمانی که رسیدیم، تا وقتی که ارتفاع ساختمان از ساختمان سر پشته بیشتر بود سر پشته را حذف می‌کنیم چرا که ساختمان سر پشته برای ساختمان‌های چپ‌تر از این ساختمان نمی‌تواند کاندید باشد. دلیل آن هم این است که اگر نوبت ساختمان \متن‌لاتین{i}ام باشد و ساختمان سر پشته ساختمان \متن‌لاتین{j}ام باشد، می‌دانیم \متن‌لاتین{h[i] > h[j]} و \متن‌لاتین{i < j} پس اگر ساختمان \متن‌لاتین{j} بخواهد اولین ساختمان بلندتر سمت راست ساختمانی مثل \متن‌لاتین{k} باشد، چون ارتفاع ساختمان \متن‌لاتین{i} بیشتر است و چپ‌تر از ساختمان \متن‌لاتین{j} هم هست، پس ساختمان \متن‌لاتین{i} هم از ساختمان \متن‌لاتین{k} بلندتر است و راست آن ساختمان است و قبل از ساختمان \متن‌لاتین{j} هم قرار دارد که تناقض است چرا که ساختمان \متن‌لاتین{i} باید اولین ساختمان بلندتر سمت راست میبوده. بنابر این عملیات حذف از سر پشته را ادامه می‌دهیم تا به یک ساختمان بلندتر برسیم. اولین جاییکه به ساختمان بلندتر از ساختمان \متن‌لاتین{i}ام در پشته برسیم، همان ساختمان \متن‌لاتین{r[i]} مورد نظر است (یعنی اولین ساختمان بلندتر از ساختمان \متن‌لاتین{i} که راست آن است) چرا که در بین کاندیدا ها زودتر از همه در پشته به آن رسیدیم یعنی هم ارتفاعش بلندتر است و هم سمت چپ کاندیداهای دیگر است پس اولین ساختمان بلندتر سمت راست ساختمان \متن‌لاتین{i}ام است. سپس ساختمان \متن‌لاتین{i} را به روی پشته اضافه می‌کنیم. 
\\کد نمونه (با تشکر از مهدی جعفری) :‌
%}
% صورت اصلی مسئله را در این‌جا بنویسید.
%اگر مسئله چند قسمت داشت به این صورت عمل کنید:
%\شروع{زیرمسائل}
%\زیرمسئله[0] % مسئلۀ ۱آ
%\زیرمسئله % مسئلۀ ۱ب
%\زیرمسئله % مسئلۀ ۱پ
%\زیرمسئله % مسئلۀ ۱ت
%\پایان{زیرمسائل}



%%%%%%%%%%%%%%%%%%%%%%%%%%%%%%%%%%%%%%%%%%%%%%%%%%%%%
% برخی دستورات مفید برای تمرین‌ها. %
%%%%%%%%%%%%%%%%%%%%%%%%%%%%%%%%%%%%%%%%%%%%%%%%%%%%%

% فرمول‌های نمونه:
%$f(x) = \Theta(x)$
%$T(x, y) \leq \log(x) + 2^y + \binom{2n}{n}$
% {\tt code\_function}


% لیست موارد بدون شماره‌گذاری:
%\شروع{فقرات}
%    \فقره مورد اول 
%        \شروع{فقرات}
%            \فقره مورد اول داخلی 
%                \شروع{فقرات}
%                    \فقره اولین مورد خیلی داخلی 
%                    \فقره دومین مورد خیلی داخلی 
%                \پایان{فقرات}
%            \فقره مورد دوم داخلی 
%        \پایان{فقرات}
%    \فقره مورد دوم 
%\پایان{فقرات}

% لیست موارد شماره‌گذاری شده:
%\شروع{شمارش}
%    \فقره یک مورد 
%    \فقره یک مورد دیگر 
%    \فقره باز هم یک مورد دیگر
%\پایان{شمارش}

% برای نوشتن یک فرمول چند خطی که یک شماره بخوره. برای حالت‌های دیگر به سایت زیر مراجعه کنید:
% https://www.overleaf.com/learn/latex/Aligning_equations_with_amsmath
%\begin{align} 
%    \begin{split}
%        (x+y)^3 &= (x+y)^2(x+y) \\
%                &= (x^2+2xy+y^2)(x+y) \\
%                &= (x^3+2x^2y+xy^2) + (x^2y+2xy^2+y^3) \\
%                &= x^3+3x^2y+3xy^2+y^3
%    \end{split}                                 
%\end{align}

% نوشتن ماتریس و آرایه دوبعدی
%\begin{align}
%    A = 
%    \begin{bmatrix}
%        A_{11} & A_{21} \\
%        A_{21} & A_{22}
%    \end{bmatrix}
%\end{align}

% اضافه کردن تصویر با شرح. برای اضافه کردن تصویر بدون شرح، فقط دستور درج‌تصویر را استفاده کنید.
%\شروع{شکل}[h!]
%\تنظیم‌ازوسط
%  \درج‌تصویر[پهنا=0.5\پهنای‌متن]{img.jpg}
%\شرح{شرح تصویر}
%\برچسب{تصویر:۱}
%\پایان{شکل}
%\پایان{زیرمسائل}

%کد زیر را ببینید.
\begin{latin}
\begin{lstlisting}
from sys import stdin, stdout


class Node:
    def __init__(self, val: int):
        self.key = val
        self.left = None
        self.right = None
        self.h = 1
        self.sum = val
        self.size = 1


tmp = []


class Tree:

    def __init__(self):
        self.root = None

    def insert(self, val: int):
        self.root = self.Insert(self.root, val)

    def Insert(self, root, val: int):
        if root is None:
            return Node(val)
        if val <= root.key:
            global tmp
            tmp[0] += Tree.getSum(root.right) + root.key
            tmp[1] += Tree.getSize(root.right) + 1
            root.left = self.Insert(root.left, val)
        else:
            root.right = self.Insert(root.right, val)

        Tree.calc(root)
        balance = Tree.getHeight(root.left) - Tree.getHeight(root.right)
        if balance > 1:
            if val <= root.left.key:
                return Tree.rightRotate(root)
            else:
                root.left = Tree.leftRotate(root.left)
                return Tree.rightRotate(root)
        elif balance < -1:
            if val <= root.right.key:
                root.right = Tree.rightRotate(root.right)
                return Tree.leftRotate(root)
            else:
                return Tree.leftRotate(root)
        else:
            return root

    @staticmethod
    def rightRotate(y):
        x = y.left
        y.left = x.right
        x.right = y

        Tree.calc(y)
        Tree.calc(x)
        return x

    @staticmethod
    def leftRotate(x):
        y = x.right
        x.right = y.left
        y.left = x

        Tree.calc(x)
        Tree.calc(y)
        return y

    @staticmethod
    def calc(node):
        node.h = max(Tree.getHeight(node.left), Tree.getHeight(node.right)) + 1
        node.size = Tree.getSize(node.left) + Tree.getSize(node.right) + 1
        node.sum = Tree.getSum(node.left) + Tree.getSum(node.right) + node.key

    @staticmethod
    def getSum(node):
        if not node:
            return 0
        return node.sum

    @staticmethod
    def getSize(node):
        if not node:
            return 0
        return node.size

    @staticmethod
    def getHeight(node):
        if not node:
            return 0
        return node.h


def main():
    global tmp
    n = int(stdin.readline())
    a = list(map(int, stdin.readline().split()))

    tree = Tree()

    ans = [0] * n
    for i in range(n - 1, -1, -1):
        tmp = [0, 0]
        tree.insert(a[i])
        ans[i] = tmp[0] - tmp[1] * a[i]

    res = [0] * n
    stack = [0] * n
    sz = 0
    for i in range(n - 1, -1, -1):
        while sz > 0 and a[stack[sz - 1]] < a[i]:
            sz -= 1
        if sz > 0:
            res[i] = ans[stack[sz - 1]]
        else:
            res[i] = 0
        stack[sz] = i
        sz += 1

    print(*res, sep=' ')


main()

\end{lstlisting}
\end{latin}

